Deux forces s'exercent sur le camion~: son poids $\vec{P}$ et la force de
freinage $\vec{F}$.

La route étant horizontale, le poids ne travaille pas.

D'après le théorème de l'énergie cinétique~:
$\Delta E_{c}=W(\vec{P})+W(\vec{F})=W(\vec{F})$

La force de freinage est constante, parallèle à la route et dirigée dans le sens
inverse du déplacement d'où~: ${W(\vec{F})=-F \times AB}$

Or~:
$\Delta E_{c}=E_{c}(B)-E_{c}(A)=\frac{1}{2}Mv_{B}^{2}-\frac{1}{2}Mv_{A}^{2}=
0-\frac{1}{2}Mv^{2}$

Donc~:
$-\frac{1}{2}Mv^{2}=-F \times AB$

Valeur de la force de freinage~:
$F=-\frac{Mv^{2}}{2\times AB}=
\frac{30\times 10^{3}\times\bigl(\frac{90\times 10^{3}}{3600}\bigr)^{2}}{2\times 500}=
\qty{18750}{\N}$\\


